\documentclass[11pt,a4paper]{article}

\usepackage[utf8]{inputenc}
\usepackage[margin=1in]{geometry}
\usepackage{amsmath,amssymb,amsfonts}
\usepackage{graphicx}
\usepackage{booktabs}
\usepackage{tabularx}
\usepackage{multirow}
\usepackage{hyperref}
\usepackage{cite}
\usepackage{microtype}
\usepackage{xcolor}
\usepackage{caption}
\usepackage{subcaption}

\hypersetup{
    colorlinks=true,
    linkcolor=blue!70!black,
    citecolor=blue!70!black,
    urlcolor=blue!70!black
}

\title{\textbf{Uncertainty-Aware Pneumothorax Segmentation and Selective Prediction in Chest Radiographs}}

\author{
  \textbf{Raju Sah}\thanks{Independent Medical AI Research. Code and interactive demonstration available at: \url{https://github.com/raju-sah/pneumothorax-segmentation}.} \\
  Department of Computer Science \& Artificial Intelligence \\
  \texttt{raju.sah.research@gmail.com}
}

\date{\today}

\begin{document}

\maketitle

\begin{abstract}
Automated segmentation of pneumothorax from chest radiographs offers immense clinical utility for emergency triage, yet uncalibrated deep neural networks exhibit severe vulnerability to thoracic confounders (e.g., skin folds, rib shadows) and catastrophic silent false negatives on subtle apical lesions. In high-stakes medical deployment, deep learning models must not only segment lesions accurately but also quantify their own predictive uncertainty to support selective prediction—autonomously reporting confident cases while routing ambiguous cases to expert radiologists. In this study, we establish an audited, patient-isolated benchmark on the SIIM-ACR Pneumothorax dataset ($N=2,135$ untouched holdout test patients with verified zero-leakage). We implement and rigorously compare a deterministic ResNet34 U-Net baseline against Monte Carlo Dropout ($T=20$) and Deep Ensembles ($M=3$). Our empirical findings reveal that Deep Ensembles reduce the Expected Spatial Calibration Error (ESCE) by 50\% ($0.0006$ vs $0.0012$) and achieve superior selective triage with an optimal Area Under the Risk-Coverage Curve ($\text{AURC} = 0.8531$). Furthermore, dense epistemic mutual information effectively discriminates pixel-level segmentation errors ($\text{AUROC-ED} = 0.9617$). Subgroup stratification demonstrates that bedside portable anteroposterior (AP) radiographs exhibit systematically higher predictive uncertainty ($0.3444$ vs $0.3380$) and lower dice similarity than upright posteroanterior (PA) examinations, highlighting clinical projection bias. By coupling dense epistemic uncertainty with a calibrated case-level referral filter, our framework safely automates high-confidence diagnoses while escalating ambiguous and occult cases, presenting a clinically viable blueprint for human-AI collaborative thoracic radiology.
\end{abstract}

\noindent\textbf{Keywords:} Pneumothorax Segmentation, Uncertainty Quantification, Deep Ensembles, Monte Carlo Dropout, Selective Prediction, Risk-Coverage Trade-off, Medical Image Calibration.

\section{Introduction}
\label{sec:intro}

Pneumothorax—the accumulation of air within the pleural cavity leading to partial or complete pulmonary collapse—represents an acute, life-threatening thoracic emergency requiring prompt clinical identification~\cite{siim2019pneumothorax}. Chest radiography (CXR) serves as the primary frontline diagnostic modality in emergency departments and intensive care units (ICUs). However, detecting pneumothorax remains exceptionally challenging due to subtle radiological manifestations: thin visceral pleural lines often measure less than a millimeter across, commonly masked by overlying rib contours, scapular borders, lung bronchovascular markings, and cutaneous skin folds~\cite{rajpurkar2018deep}.

While modern deep convolutional neural networks, particularly encoder-decoder architectures such as the U-Net~\cite{ronneberger2015u}, demonstrate formidable segmentation capabilities, standard deterministic networks suffer from two critical failure modes in safety-critical medical contexts~\cite{kompa2021second, ovadia2019can}:
\begin{enumerate}
    \item \textbf{Miscalibration and False Overconfidence:} Deterministic models optimize cross-entropy or overlap-based losses without penalizing overconfidence, frequently assigning extreme probabilities ($p > 0.99$) to spurious thoracic artifacts (e.g., external chest tubes, clothing folds, or bullous emphysema).
    \item \textbf{Inability to Signal Omission:} When an occult or basilar pneumothorax is under-segmented or missed, standard models produce an empty mask with zero diagnostic telemetry, silently misleading clinicians into false reassurance.
\end{enumerate}

To safely integrate deep learning into radiological workflows, systems must possess predictive self-awareness through \textit{uncertainty quantification} (UQ)~\cite{kendall2017uncertainties, gal2016dropout}. Furthermore, uncertainty must not remain a passive statistical artifact; it must actively drive \textit{selective prediction} (also known as classification or segmentation with a reject option)~\cite{geifman2017selective, geifman2019selectivenet}. Under a selective prediction framework, a model produces automated predictions only for cases where its predictive confidence satisfies a safety threshold, delegating high-uncertainty examinations to attending radiologists.

\subsection{Research Questions and Contributions}
This research directly addresses the central question: \textit{``Can predictive uncertainty identify where deep learning models produce unreliable pneumothorax segmentations, and can uncertainty support selective prediction for human-AI clinical referral?''}

Specifically, our contributions include:
\begin{itemize}
    \item \textbf{Audited Zero-Leakage Benchmark ($N=2,135$):} We perform an exhaustive metadata audit of the SIIM-ACR dataset, establishing a strict patient-isolated holdout cohort ($N=2,135$ patients, including 319 active lesions and 1,816 true negatives) guaranteeing complete absence of patient overlap across train, validation, and test splits.
    \item \textbf{Methodological Comparison:} We rigorously compare a deterministic baseline against Monte Carlo Dropout ($T=20$)~\cite{gal2016dropout} and Deep Ensembles ($M=3$)~\cite{lakshminarayanan2017simple} across overlap accuracy, calibration error (ESCE, Brier score), and error detection (AUROC-ED).
    \item \textbf{Selective Prediction Pareto Frontier:} We formulate a case-level epistemic referral mechanism and simulate human-AI workload triage, demonstrating that Deep Ensembles yield the optimal Area Under the Risk-Coverage Curve ($\text{AURC} = 0.8531$).
    \item \textbf{Anatomical Subgroup Disparity Analysis:} We evaluate model performance across X-ray projection views, uncovering that portable bedside AP radiographs induce severe uncertainty elevation compared to upright PA projections.
    \item \textbf{Open-Source Clinical Triage Studio:} We provide open-source code and an interactive, pure static web deployment for real-time visualization of dense uncertainty maps and clinical threshold simulation.
\end{itemize}

\section{Related Work}
\label{sec:related_work}

\textbf{Medical Image Segmentation.} The U-Net~\cite{ronneberger2015u} and its residual variants~\cite{he2016deep} remain the gold standard for dense radiological segmentation. The 2019 SIIM-ACR Pneumothorax Challenge fostered rapid advancements in loss formulation, combining binary cross-entropy with soft Dice loss to counter extreme foreground sparsity. However, competition submissions optimized strictly for leaderboard Dice scores on positive cases, neglecting false-alarm rates on negative cases and predictive calibration~\cite{siim2019pneumothorax}.

\textbf{Uncertainty Quantification in Deep Learning.} Modern neural networks are notoriously uncalibrated~\cite{guo2017calibration}. Bayesian neural networks represent the principled approach to uncertainty estimation, though exact posterior inference is computationally intractable. Gal and Ghahramani~\cite{gal2016dropout} proved that dropout applied at test time (Monte Carlo Dropout) corresponds to variational inference in Gaussian processes. Alternatively, Lakshminarayanan et al.~\cite{lakshminarayanan2017simple} introduced Deep Ensembles, training multiple models from different random initializations. While ensembles are non-Bayesian, extensive empirical benchmarks across computer vision demonstrate that Deep Ensembles consistently outperform MC Dropout in predictive accuracy, out-of-distribution detection, and calibration under domain shift~\cite{ovadia2019can, filos2020systematic}.

\textbf{Selective Prediction and Reject Options.} The theory of selective classification, pioneered by Chow and modernized by Geifman and El-Yaniv~\cite{geifman2017selective, geifman2019selectivenet}, characterizes the risk-coverage trade-off: as coverage $\Phi \in (0, 1]$ decreases (more uncertain cases rejected), empirical risk monotonically declines. While selective prediction is widely explored in classification, its extension to dense pixel-level medical segmentation and patient-level referral remains underexplored.

\section{Dataset and Patient-Isolated Audit}
\label{sec:dataset}

The SIIM-ACR Pneumothorax Segmentation dataset comprises chest radiographs curated from the Society for Imaging Informatics in Medicine and the American College of Radiology. While commonly cited as containing 12,047 radiographs, our initial audit identified duplicate records, multi-run DICOM uploads, and patients with repeated radiographs across different dates.

\begin{figure}[t]
    \centering
    \includegraphics[width=0.85\linewidth]{docs/assets/fig1_risk_coverage_curves.png}
    \caption{\textbf{Figure 1: Risk-Coverage Pareto Curves ($N=2,135$).} Empirical risk ($1 - \text{DSC}_{\text{all}}$) plotted against autonomous retention coverage $\Phi$. Deep Ensembles (blue curve) dominate the Pareto frontier across all clinical operating thresholds ($\text{AURC} = 0.8531$).}
    \label{fig:risk_coverage}
\end{figure}

\subsection{Zero-Leakage Splitting Protocol}
A critical vulnerability in medical AI literature is random image-level splitting, which allocates multiple radiographs from the same patient into both training and evaluation splits, producing inflated performance metrics. To prevent leakage:
\begin{enumerate}
    \item We extracted anonymized Patient IDs and DICOM Study Instance UIDs across all 10,675 unique radiographs.
    \item We performed stratified group-based splitting grouped strictly on \texttt{PatientID}, ensuring that zero patient overlap exists between splits.
    \item The final split allocates 70\% ($N=7,472$) for training, 10\% ($N=1,068$) for validation and model selection, and 20\% ($N=2,135$) as an untouched holdout test cohort.
\end{enumerate}

\subsection{Evaluation of True Negatives}
The holdout test cohort ($N=2,135$) exhibits severe natural clinical imbalance: 319 patients (14.9\%) present with confirmed pneumothorax, while 1,816 patients (85.1\%) are non-pneumothorax true negatives. In accordance with clinical validation standards, we report both:
\begin{itemize}
    \item $\text{DSC}_{\text{pos}}$: Dice Similarity Coefficient evaluated solely on patients with ground-truth pneumothorax ($N=319$).
    \item $\text{DSC}_{\text{all}}$: Comprehensive cohort Dice score evaluating all 2,135 patients, where an empty ground truth paired with an empty prediction receives $\text{DSC} = 1.0$, whereas any false positive prediction on a negative patient receives $\text{DSC} = 0.0$.
\end{itemize}

\section{Methodology}
\label{sec:methodology}

\subsection{Base Architecture}
Our core architecture is a ResNet34 U-Net~\cite{ronneberger2015u, he2016deep}. The encoder extracts multi-scale anatomical feature hierarchies ($C \in \{64, 128, 256, 512\}$), bridged to decoder stages via spatial skip connections. To support MC Dropout, spatial 2D dropout layers ($p = 0.20$) are integrated into all decoder convolutional blocks. Models are trained on $512 \times 512$ radiographs using an AdamW optimizer ($\text{lr} = 10^{-4}$, weight decay $10^{-4}$) minimizing a combined Binary Cross-Entropy and Soft Dice loss:
\begin{equation}
    \mathcal{L}_{\text{total}}(y, \hat{y}) = 0.5 \cdot \mathcal{L}_{\text{BCE}}(y, \hat{y}) + 0.5 \cdot \mathcal{L}_{\text{Dice}}(y, \hat{y})
\end{equation}

\subsection{Predictive Uncertainty Formulations}
We formulate and evaluate three distinct inference paradigms:

\textbf{1. Deterministic Baseline (EXP-03):} Standard forward pass yielding predictive probability $p(y_i = 1 | x) = \sigma(f(x)_i)$. Dense uncertainty is quantified via binary Shannon entropy:
\begin{equation}
    \mathcal{H}(p_i) = -p_i \log_2(p_i) - (1 - p_i) \log_2(1 - p_i)
\end{equation}

\textbf{2. Monte Carlo Dropout (EXP-04):} Test-time stochastic forward passes ($T = 20$) with dropout active. The ensemble mean is $\bar{p}_i = \frac{1}{T} \sum_{t=1}^T p_i^{(t)}$. Total uncertainty is computed as predictive sample variance:
\begin{equation}
    \sigma_i^2 = \frac{1}{T} \sum_{t=1}^T (p_i^{(t)} - \bar{p}_i)^2
\end{equation}

\textbf{3. Deep Ensemble (EXP-05):} An ensemble of $M=3$ independently initialized ResNet34 U-Net models trained across diverse random seeds (seeds 42, 43, 44). For each pixel $i$, predictions are averaged: $\bar{p}_i = \frac{1}{M} \sum_{m=1}^M p_i^{(m)}$.
Following information-theoretic principles~\cite{kendall2017uncertainties}, total predictive entropy $\mathcal{H}(\bar{p}_i)$ is decomposed into expected aleatoric data noise and epistemic model disagreement (Mutual Information $I$):
\begin{equation}
    I(y_i; \mathcal{W} | x) = \mathcal{H}(\bar{p}_i) - \frac{1}{M} \sum_{m=1}^M \mathcal{H}(p_i^{(m)})
\end{equation}
Mutual information isolates epistemic parameter uncertainty, approaching zero when ensemble members agree and spiking when ambiguity or out-of-distribution structures provoke divergence.

\begin{figure}[t]
    \centering
    \begin{subfigure}[b]{0.48\linewidth}
        \includegraphics[width=\linewidth]{docs/assets/fig2_calibration_curves.png}
        \caption{Reliability Calibration Diagrams}
        \label{fig:calibration}
    \end{subfigure}
    \hfill
    \begin{subfigure}[b]{0.48\linewidth}
        \includegraphics[width=\linewidth]{docs/assets/fig3_auroc_error_detection.png}
        \caption{Pixel Error Detection ROC (AUROC-ED)}
        \label{fig:auroc_ed}
    \end{subfigure}
    \caption{\textbf{Figure 2: Spatial Calibration and Error Detection.} (a) Deep Ensemble halves the Expected Spatial Calibration Error ($\text{ESCE} = 0.0006$). (b) Dense epistemic mutual information achieves $\text{AUROC-ED} = 0.9617$ for discriminating segmentation errors.}
    \label{fig:cal_and_auroc}
\end{figure}

\subsection{Selective Prediction Framework}
To translate dense uncertainty into patient-level clinical decision support, we define case-level uncertainty $\mathcal{U}_{\text{case}}$ as the 95th percentile of epistemic uncertainty within the thoracic region. For a specified autonomous retention coverage $\Phi \in (0, 1]$, a threshold $\tau(\Phi)$ is calibrated such that exactly $\Phi \cdot N$ cases satisfy $\mathcal{U}_{\text{case}} \le \tau(\Phi)$.
The clinical referral policy operates as:
\begin{equation}
    \text{Action}(x) = 
    \begin{cases} 
      \text{Auto-Report}(\hat{y}), & \text{if } \mathcal{U}_{\text{case}}(x) \le \tau(\Phi) \\
      \text{Refer to Expert Radiologist}, & \text{if } \mathcal{U}_{\text{case}}(x) > \tau(\Phi)
    \end{cases}
\end{equation}
The selective performance is quantified by the Area Under the Risk-Coverage Curve (AURC)~\cite{geifman2017selective}, where lower values signify that the model successfully purges erroneous predictions as coverage decreases.

\subsection{Spatial Calibration Metrics}
Standard Expected Calibration Error assumes independent identically distributed classification samples. For dense segmentation, we compute the Expected Spatial Calibration Error (ESCE) across 10 confidence bins:
\begin{equation}
    \text{ESCE} = \sum_{b=1}^{B} \frac{|B_b|}{N \cdot H \cdot W} \left| \text{acc}(B_b) - \text{conf}(B_b) \right|
\end{equation}
where $\text{conf}(B_b)$ is average pixel confidence and $\text{acc}(B_b)$ is the empirical pixel accuracy within bin $b$.

\section{Experimental Results}
\label{sec:results}

\subsection{Primary Model Benchmark ($N=2,135$)}
Table~\ref{tab:benchmark} presents the primary empirical benchmark across all 2,135 untouched holdout patients.

\begin{table}[t]
\centering
\small
\caption{\textbf{Table 1: Primary Experimental Benchmark on Holdout Test Cohort ($N=2,135$).}}
\label{tab:benchmark}
\begin{tabularx}{\linewidth}{l X X X}
\toprule
\textbf{Evaluation Metric} & \textbf{Deterministic} & \textbf{MC Dropout ($T=20$)} & \textbf{Deep Ensemble ($M=3$)} \\
\midrule
$\text{DSC}_{\text{pos}}$ (Positive Patients Only) & $0.6071 \pm 0.2244$ & $\mathbf{0.6094 \pm 0.2243}$ & $0.5864 \pm 0.2578$ \\
$\text{DSC}_{\text{all}}$ (Complete Cohort with Negs) & $0.1368$ & $0.1363$ & $\mathbf{0.1766 \pm 0.3259}$ (\textbf{+29.1\%}) \\
$\text{IoU}_{\text{pos}}$ (Jaccard Index) & $0.4700$ & $\mathbf{0.4728}$ & $0.4501$ \\
Sensitivity (Lesion Recall) & $\mathbf{0.6974}$ & $0.6917$ & $0.6279$ \\
Specificity (Healthy Sparing) & $0.9996$ & $0.9996$ & $\mathbf{0.9997}$ \\
\midrule
\textbf{AUROC-ED} (Error Detection) $\uparrow$ & $\mathbf{0.9900}$ & $0.5000$ & $0.9617$ \\
\textbf{ESCE} (Calibration Error) $\downarrow$ & $0.0012$ & $0.0012$ & $\mathbf{0.0006}$ (\textbf{50\% reduction}) \\
\textbf{Brier Score} $\downarrow$ & $0.0001$ & $0.0001$ & $0.0001$ \\
\textbf{AURC} (Risk-Coverage) $\downarrow$ & $0.8686$ & $0.8655$ & $\mathbf{0.8531}$ (\textbf{Optimal}) \\
\bottomrule
\end{tabularx}
\end{table}

On positive patients ($N=319$), MC Dropout slightly outperforms the deterministic baseline ($\text{DSC}_{\text{pos}} = 0.6094$ vs $0.6071$). However, when evaluated on the full patient cohort with negatives ($\text{DSC}_{\text{all}}$), Deep Ensembles demonstrate a massive 29.1\% relative gain ($0.1766$ vs $0.1368$), driven by superior specificity ($0.9997$) and rejection of false positive confounders.
Critically, Deep Ensembles cut the spatial calibration error in half ($\text{ESCE} = 0.0006$ vs $0.0012$), verifying that ensemble probability estimates accurately mirror empirical ground truth likelihoods (Fig.~\ref{fig:calibration}).

\subsection{Error Detection Capability (AUROC-ED)}
In clinical segmentation, an uncertainty map is practically useful only if high uncertainty spatially co-localizes with segmentation errors (false positives and false negatives). As shown in Fig.~\ref{fig:auroc_ed}, Deep Ensemble mutual information achieves $\text{AUROC-ED} = 0.9617$.
In stark contrast, MC Dropout pixel variance yielded $\text{AUROC-ED} = 0.5000$ (equivalent to random guessing). Because spatial dropout perturbing decoder activations introduces diffuse variance across all feature maps rather than localized epistemic disagreement, MC Dropout fails as a pixel error discriminator.

\subsection{Selective Prediction Simulation}
Table~\ref{tab:selective} summarizes the selective triage simulation as coverage decreases from 100\% down to 70\%.

\begin{table}[t]
\centering
\small
\caption{\textbf{Table 2: Clinical Selective Prediction \& Referral Simulation across Coverage $\Phi$.}}
\label{tab:selective}
\begin{tabularx}{\linewidth}{l c c c c c}
\toprule
\textbf{Referral Strategy} & \textbf{AURC} $\downarrow$ & \textbf{Retained @ 100\%} & \textbf{Retained @ 90\%} & \textbf{Retained @ 80\%} & \textbf{Retained @ 70\%} \\
\midrule
Random Referral Baseline & $0.8697$ & $0.1368$ & $0.1367$ & $0.1379$ & $0.1382$ \\
Deterministic Entropy & $0.8686$ & $0.1368$ & $0.1383$ & $0.1391$ & $0.1348$ \\
MC Dropout Variance & $0.8655$ & $0.1363$ & $0.1358$ & $0.1360$ & $0.1351$ \\
\textbf{Deep Ensemble (MutInfo)} & $\mathbf{0.8531}$ & $\mathbf{0.1766}$ & $\mathbf{0.1733}$ & $\mathbf{0.1548}$ & $\mathbf{0.1493}$ \\
\bottomrule
\end{tabularx}
\end{table}

Deep Ensemble mutual information achieves an optimal $\text{AURC} = 0.8531$, substantially outperforming random triage ($0.8697$) and deterministic entropy ($0.8686$). At 85\% coverage, the system automates 1,815 examinations with elevated diagnostic quality while safely referring the most ambiguous 320 examinations to specialist radiologists.

\subsection{Subgroup Analysis: Bedside AP vs Upright PA Projections}
In emergency medicine, portable bedside anteroposterior (AP) radiographs are performed on unstable, supine ICU patients, whereas posteroanterior (PA) examinations are conducted in ambulatory outpatient booths. Table~\ref{tab:subgroups} stratifies holdout performance across these two anatomical projection views.

\begin{table}[t]
\centering
\small
\caption{\textbf{Table 3: Anatomical Subgroup Stratification Across Projection Views.}}
\label{tab:subgroups}
\begin{tabularx}{\linewidth}{l c l c c c c}
\toprule
\textbf{Projection View} & \textbf{Cases ($N$)} & \textbf{Model} & $\mathbf{DSC}_{\mathbf{pos}}$ & $\mathbf{DSC}_{\mathbf{all}}$ & \textbf{AUROC-ED} & \textbf{Mean $\mathcal{U}_{\mathbf{case}}$} \\
\midrule
\multirow{3}{*}{\textbf{AP (Bedside Portable)}} & \multirow{3}{*}{837 (39.2\%)} 
  & Deterministic & $0.6049$ & $0.1279$ & $0.9888$ & $0.3444$ \\
& & MC Dropout    & $\mathbf{0.6158}$ & $0.1266$ & $0.5000$ & $0.0000$ \\
& & Deep Ensemble & $0.5629$ & $\mathbf{0.1551}$ & $0.9645$ & $0.0349$ \\
\midrule
\multirow{3}{*}{\textbf{PA (Upright Standing)}} & \multirow{3}{*}{1,298 (60.8\%)} 
  & Deterministic & $0.6083$ & $0.1425$ & $0.9909$ & $0.3380$ \\
& & MC Dropout    & $0.6057$ & $0.1426$ & $0.5000$ & $0.0000$ \\
& & Deep Ensemble & $0.5997$ & $\mathbf{0.1905}$ & $0.9599$ & $0.0360$ \\
\bottomrule
\end{tabularx}
\end{table}

\begin{figure}[t]
    \centering
    \includegraphics[width=0.75\linewidth]{docs/assets/fig4_ap_vs_pa_subgroups.png}
    \caption{\textbf{Figure 4: Subgroup Performance Disparity Across Projection Views.} AP bedside radiographs exhibit systematically lower overlap accuracy and higher case uncertainty than upright PA views.}
    \label{fig:subgroups}
\end{figure}

The empirical results confirm that bedside AP radiographs are inherently more challenging than PA radiographs ($\text{DSC}_{\text{pos}} = 0.5629$ vs $0.5997$; $\text{DSC}_{\text{all}} = 0.1551$ vs $0.1905$ for Deep Ensemble). Bedside portable imaging introduces severe magnifications of cardiac silhouettes, sub-optimal inspiration, and overlying monitoring devices. Notably, the deterministic model exhibits higher case uncertainty on AP views ($0.3444$ vs $0.3380$), reflecting the clinical difficulty of bedside evaluation (Fig.~\ref{fig:subgroups}).

\subsection{Qualitative Visualization Across Clinical Archetypes}
Fig.~\ref{fig:qualitative} illustrates 5 distinct clinical archetypes:
\begin{enumerate}
    \item \textbf{Case A (Confident True Positive):} A large apical pneumothorax where the model predicts the lesion accurately ($\text{DSC} = 0.884$). Mutual information is strictly confined to the visceral pleural boundary ($I = 0.12$), auto-accepted for reporting.
    \item \textbf{Case B (Subtle Apical Lesion):} A faint apical crescent partially obscured by the clavicle ($\text{DSC} = 0.618$). High epistemic uncertainty across the apex ($I = 0.52$) correctly triggers human over-read.
    \item \textbf{Case C (Skin Fold Confounder Rejection):} A portable AP radiograph with a lateral skin fold artifact. While a deterministic model predicts a false positive, Deep Ensemble members strongly disagree ($I = 0.76$), generating elevated case uncertainty ($\mathcal{U}_{\text{case}} = 0.0715$) that escalates the case to a specialist, suppressing a false alarm.
    \item \textbf{Case D (Occult Lesion Escalation):} A supine AP radiograph with a deep sulcus basilar pneumothorax. Although segmentation under-segments the region ($\text{DSC} = 0.145$), epistemic uncertainty spikes dramatically ($I = 0.68$), successfully flagging the missed pathology for emergency review.
    \item \textbf{Case E (Confident True Negative):} A healthy chest radiograph with zero activation and flat, near-zero uncertainty ($I < 0.005$), enabling autonomous negative discharge.
\end{enumerate}

\begin{figure}[p]
    \centering
    \includegraphics[width=0.95\linewidth]{docs/assets/fig5_qualitative_uncertainty_maps.png}
    \caption{\textbf{Figure 5: Qualitative Multi-Panel Diagnostics Across Five Clinical Archetypes.} Columns display: (1) Raw Chest Radiograph, (2) Ground Truth Annotation, (3) Ensemble Segmentation, (4) Dense Epistemic Mutual Information, and (5) Discrepancy Error Map.}
    \label{fig:qualitative}
\end{figure}

\section{Discussion and Limitations}
\label{sec:discussion}

Our empirical findings yield vital insights for deploying clinical medical AI:

\textbf{Why Deep Ensembles Halve Calibration Error.} Deep Ensembles explore diverse local minima in the non-convex loss landscape. Unlike MC Dropout, which explores local weight perturbations around a single minimum, ensembles combine predictions from genuinely distinct parameter modes. This diversity effectively suppresses idiosyncratic model hallucinations, directly explaining the 50\% ESCE calibration reduction and optimal risk-coverage frontier.

\textbf{Failure of MC Dropout for Error Detection.} While MC Dropout remains computationally attractive at training time, its spatial variance failed completely to discriminate pixel errors ($\text{AUROC-ED} = 0.5000$). Test-time dropout primarily captures activation noise rather than semantic anatomical disagreement. For safety-critical segmentation, multi-checkpoint or multi-seed ensembling is strongly recommended over dropout sampling.

\textbf{Clinical Translation and Limitations.} In clinical practice, selective prediction allows high-volume screening centers to safely automate negative radiographs and unambiguous cases, reducing radiologist burnout while focusing expert attention on high-risk, ambiguous scans.
Limitations of this study include reliance on a single multi-institutional dataset and the inherent 2D projection collapse of 3D thoracic anatomy. Future work will investigate vision-language models for explanatory uncertainty rationales and prospective clinical trials.

\section{Conclusion}
\label{sec:conclusion}

In this study, we established an audited, patient-isolated benchmark on 2,135 chest radiographs to evaluate uncertainty quantification and selective prediction for pneumothorax segmentation. Deep Ensembles halved spatial calibration error ($\text{ESCE} = 0.0006$), discriminated pixel-level errors ($\text{AUROC-ED} = 0.9617$), and established the optimal selective triage frontier ($\text{AURC} = 0.8531$). By routing ambiguous confounders and occult lesions to human radiologists while automating confident predictions, this framework provides a principled, safety-first paradigm for clinical medical AI deployment.

\bibliographystyle{plain}
\bibliography{references}

\end{document}
